% Source: https://artofproblemsolving.com/community/c3h1990481_puzzleriddle_extended
% Stable identity: aops-topic-1990481-q01
Ask either guard:
\begin{quote}
"If I asked you whether the left door leads to the prize, would you answer \emph{Ghi}?"
\end{quote}
\textbf{Claim:} the guard answers "Ghi" if and only if the left door leads to the prize. So if he says "Ghi", choose the left door; if he says "Pon", choose the right door.

\textbf{Proof.} Let $L$ be the statement "the left door leads to the prize". There are four cases.
\begin{itemize}
\item \emph{Truth-teller, Ghi means yes.} Asked about $L$, he would say Ghi exactly when $L$ is true. So the asked statement is true iff $L$, and he truthfully answers yes (Ghi) iff $L$.
\item \emph{Truth-teller, Ghi means no.} Asked about $L$, he would say Ghi (no) exactly when $L$ is false, so the asked statement is true iff not $L$. Answering truthfully, he says yes (Pon) iff not $L$, i.e.\ he says Ghi iff $L$.
\item \emph{Liar, Ghi means yes.} Asked about $L$, he would lie and say Ghi (yes) exactly when $L$ is false, so the statement is true iff not $L$. He lies about it, saying yes (Ghi) iff the statement is false, i.e.\ iff $L$.
\item \emph{Liar, Ghi means no.} Asked about $L$, he would lie and say Ghi (no) exactly when $L$ is true, so the statement is true iff $L$. He lies, saying yes (Pon) iff the statement is false, i.e.\ iff not $L$; so he says Ghi iff $L$.
\end{itemize}
In every case the answer is "Ghi" exactly when the left door leads to the prize. $\blacksquare$
